%---------------------------------------------------------------------------%
%->> Backmatter
%---------------------------------------------------------------------------%
\chapter[致谢]{致\quad 谢}\chaptermark{致\quad 谢}% syntax: \chapter[目录]{标题}\chaptermark{页眉}
%\thispagestyle{noheaderstyle}% 如果需要移除当前页的页眉
%\pagestyle{noheaderstyle}% 如果需要移除整章的页眉

光阴荏苒，白驹过隙。回首在中科院信工所求学的这几年，从最初敲下第一行代码的懵懂，到如今这本厚厚论文的落笔，心中充满了无尽的感慨与感激。这不仅是一段探索未知、叩问学术的旅程，更是一段不断褪去青涩、重塑自我的岁月。在此，我谨向所有在这个过程中给予我指导、支持与陪伴的人，致以最诚挚的谢意。

首先，感谢我的导师戴娇老师。感谢您在研究生期间为我提供了良好的科研平台与充足的研究资源，让我得以在一个包容、严谨的学术环境中顺利开展各项工作。感谢您在论文开题、中期考核及答辩等各个环节中给予的指导与建议，为我硕士学业的顺利完成提供了坚实的保障。

其次，我要向韩冀中老师致以最深切的敬意与感谢。在我的求学生涯中，您实实在在地行使着传道、授业、解惑的恩师职责。在科研上，是您高屋建瓴的指导和敏锐的学术洞察力，为我拨开了无数次研究方向上的迷雾；在生活中，您更是给予了我无微不至的关怀与点拨。您的严谨治学与豁达为人，将是我未来长久受用的宝贵财富。

我还要特别感谢北京航空航天大学的刘偲教授。从本科阶段的科研启蒙，到硕士期间的持续引领，您是我学术道路上当之无愧的引路人。是您多年如一日的悉心栽培和包容鼓励，让我从一个门外汉，逐步建立起对人工智能领域的认知与热爱。漫漫长路，师恩如海，学生铭记于心。

行文至此，我最想深深感谢的，是我的廖越师兄。在科研那些最难熬的泥泞里，是你手把手地带我推演逻辑、修改代码、复盘一次次的实验失败。你给予我的，远不止学术上毫无保留的倾囊相授，更是生活中的莫逆之交。在我陷入迷茫或内耗时，是你一次次将我稳稳接住，教会我如何去面对挫折，如何去拆解人生的难题，最重要的是，你教会了我如何去成为一个更好、更坚韧、更纯粹的人。这篇论文的每一个字符背后，都藏着你倾注的心血与耐心。

此外，还要感谢实验室的所有同门与挚友。感谢我们在无数个熬夜跑实验的日子里相互打气，在那些嬉笑怒骂的日常里分享彼此的喜怒哀乐。是你们的存在，让这段原本枯燥的代码岁月变得生动且鲜活。

最后，感谢我的家人。无论我身在何方，面临何种抉择，你们永远是我身后最坚实的后盾和最温暖的港湾。你们的默默付出与无条件的支持，是我不断前行的最大动力。

长路漫漫，过往皆为序章。写下这段文字，意味着中科院的这段旅程即将到站。愿在未来的日子里，我能带着这份沉甸甸的期许与感恩，听凭风引，继续去探索更大的世界。


\rightline{2026年4月}
\chapter{作者简历及攻读学位期间发表的学术论文与其他相关学术成果}

\section*{作者简历：}
2019年09月——2023年06月，在北京航空航天大学获得工学学士（人工智能）与理学学士（数学与应用数学）双学位

2023年09月——至今，在中国科学院信息工程研究所攻读网络空间安全硕士学位。

\section*{已发表（或正式接受）的学术论文：}

{
\setlist[enumerate]{}% restore default behavior
\begin{enumerate}[nosep]
    \item Latent DPO for Concept Erasure in Text-to-Video Diffusion Models (第一作者，已被 ICASSP 2026 接收) 
\end{enumerate}
}

\section*{参加的研究项目及获奖情况：}


\cleardoublepage[plain]% 让文档总是结束于偶数页，可根据需要设定页眉页脚样式，如 [noheaderstyle]
%---------------------------------------------------------------------------%
